\section{Heisenberg representation of operators}

\SourcePage{57}{60}
In the mathematical apparatus of quantum mechanics as here presented, the
operators corresponding to the various physical quantities act on functions
of the coordinates and as such ordinarily contain no explicit dependence on
time. The dependence of the mean values of physical quantities on time arises
only through the time dependence of the wave function of the state, in
accordance with the formula
\begin{equation}
  \bar f(t)=\int\Psi^*(q,t)\hat f\Psi(q,t)\,dq.
  \tag{13.1}
\end{equation}

The apparatus of quantum mechanics can, however, also be formulated in a
somewhat different, equivalent form, in which the dependence on time is
transferred from the wave functions to the operators. Although in this book
we shall not make use of such a representation of operators---the so-called
Heisenberg representation, as distinct from the Schr\"odinger
representation---we shall formulate it here, with later applications in
relativistic theory in view.

Let us introduce the unitary operator (cf.\ (12.13))
\begin{equation}
  \hat S=\exp\left(-\frac{i}{\hbar}\hat Ht\right),
  \tag{13.2}
\end{equation}
where $\hat H$ is the Hamiltonian of the system. By definition its
eigenfunctions coincide with the eigenfunctions of the operator $\hat H$,
i.e.\ with the wave functions $\psi_n(q)$ of the stationary states, and
moreover
\begin{equation}
  \hat S\psi_n(q)
  =\exp\left(-\frac{i}{\hbar}E_nt\right)\psi_n(q)=\Psi_n(q,t).
  \tag{13.3}
\end{equation}
It follows that expansion (10.3) of any wave function
$\Psi$ over the stationary-state wave functions can be recast in
operator form as
\begin{equation}
  \Psi(q,t)=\hat S\Psi(q,0),
  \tag{13.4}
\end{equation}
i.e.\ the action of the operator $\hat S$ carries the wave function of the
system at some initial instant of time over into the wave function at an
arbitrary instant of time.
\SourcePage{58}{61}
On introducing, in accordance with (12.7), the time-dependent operator
\begin{equation}
  \hat f(t)=\hat S^{-1}\hat f\hat S,
  \tag{13.5}
\end{equation}
we obtain
\begin{equation}
  \bar f(t)=\int\Psi^*(q,0)\hat f(t)\Psi(q,0)\,dq,
  \tag{13.6}
\end{equation}
i.e.\ we bring the formula for the mean value of the quantity $f$ (which
serves as the definition of the operators) into a form in which the
dependence on time is wholly transferred to the operator.

It is obvious that the matrix elements of the operator (13.5), with respect to
the wave functions of stationary states, coincide with the time-dependent
matrix elements $f_{nm}(t)$ defined by formula (11.3).

Finally, on differentiating expression (13.5) with respect to time (assuming
here that the operators $\hat f$ and $\hat H$ themselves contain no $t$), we
obtain the equation
\begin{equation}
  \frac{\partial\hat f(t)}{\partial t}
  =\frac{i}{\hbar}\left[\hat H\hat f(t)-\hat f(t)\hat H\right],
  \tag{13.7}
\end{equation}
which is analogous to formula (9.2) but has a somewhat different meaning:
expression (9.2) constitutes the definition of the operator $\hat{\dot f}$
corresponding to the physical quantity $\dot f$, whereas the left-hand side of
equation (13.7) contains the time derivative of the operator of the quantity
$f$ itself.
